% !TeX root = courtade-kumar.tex
% !TEX TS-program = pdfLaTeX
\section{The induction}\label{sec:induction}
All information quantities are measured in bits. For real $a,b\in[0,1]$,
write
\[
 \kldivb{a}{b}
 =a\log_2\frac ab+(1-a)\log_2\frac{1-a}{1-b},
\]
with the usual limiting values at the endpoints. For a bit $B$ of mean
$\mu$, mutual information is the average posterior divergence:
\begin{equation}\label{eq:posterior-divergence}
 \muti{B}{Z}=\E\sparen{\kldivb{\Pr(B=1\mid Z)}{\mu}}.
\end{equation}
Fix $\rho\in[0,1]$ and put $p=(1-\rho)/2$.

\subsection{The information profile}\label{sec:proof}
The divergence from a fair bit will occur repeatedly. Write
\begin{equation}\label{eq:app-entropyseries}
 \kappa(t):=\kldivb{(1+t)/2}{1/2}
 =\sum_{k\ge1}\frac{t^{2k}}{(\ln2)\,2k(2k-1)},
 \qquad \frac{t^2}{2\ln2}\le\kappa(t)\le t^2.
\end{equation}
The expansion holds also at $|t|=1$, where its coefficients sum to one.
In particular, $\muti{Y_1}{X}=\kldivb{p}{1/2}=\kappa(\rho)$.
With $\nu=4\mu(1-\mu)$, the profile in \eqref{eq:q} is
\begin{equation}\label{eq:q-divergence}
 Q_\rho(\mu)=\rho^2\nparen{1-\nu-\kldivb{\mu}{1/2}}
 +\nu\kappa(\rho).
\end{equation}

\begin{lemma}\label{lem:q}
For every $\rho,\mu\in[0,1]$,
$Q_\rho(\mu)\le\muti{Y_1}{X}$.
\end{lemma}
\begin{proof}
By \eqref{eq:app-entropyseries},
\begin{align*}
 \muti{Y_1}{X}-Q_\rho(\mu)
 &=\rho^2\kldivb{\mu}{1/2}
   -(1-\nu)\nparen{\rho^2-\kldivb{p}{1/2}}\\
 &\ge\nparen{\frac1{\ln2}-1}\rho^2(1-\nu)\ge0.
\end{align*}
\end{proof}

\subsection{Induction on the coordinates}
We prove \autoref{thm:ck} by induction on the number of coordinates,
retaining the mean of the output in $Q_\rho$. The induction uses a
coordinate whose two restrictions have sufficiently close means.

\begin{lemma}\label{lem:coordinate-gap}
Let $B=f(Y)$, where $Y$ is uniform on $\cube$, and let
$\mu=\E[B]$ and $\nu=4\mu(1-\mu)$. Among any three distinct
coordinates there is an $i$ such that
\[
 d_i:=\abs{\E[B\emid Y_i=1]-\E[B\emid Y_i=0]}
 \le \frac{\sqrt\nu}{2}.
\]
\end{lemma}
\begin{proof}
Complement $B$ if necessary so that $\mu\ge1/2$, and write $F=2B-1$.
For the three coordinates, choose signs so that
$T$ is the sum of their signed bits $\pm(2Y_i-1)$ and
$\E[FT]$ is the sum of their absolute correlations with $F$.
Since $\E|T+1|=3/2$ and $\E T=0$, the inequalities
$F(T+1)\le|T+1|$ and $F(T+3)\le T+3$ give
\[
 \E[FT]\le\min\{5/2-2\mu,\,6(1-\mu)\}
 \le\frac32\sqrt\nu.
\]
For the last step, square the first bound when $\mu\le7/8$
and the second when $\mu\ge7/8$; the respective differences are
$(2\mu-1)(25-26\mu)/4$ and $9(1-\mu)(5\mu-4)$, both nonnegative.
The smallest correlation is at most one third of their sum.
As $d_i=|\E[F(2Y_i-1)]|$, this proves the claim.
\end{proof}

The following lemma combines the information bounds for the two
restrictions. Its proof is given in \autoref{app:closure}.

\begin{lemma}\label{lem:closure}
Let $S$ be a uniform bit independent of a finite random variable
$W$, and let $B$ be a bit with any conditional law given $(S,W)$.
Obtain $S_\rho$ by flipping $S$ with probability $p$, independently
of $(S,W,B)$. Set
\[
 \mu_s=\E[B\emid S=s],\qquad
 \mu=\frac{\mu_0+\mu_1}{2},\qquad \nu=4\mu(1-\mu),\qquad d=|\mu_1-\mu_0|.
\]
If $4d^2\le\nu$ and
\begin{equation}\label{eq:child-information}
 \muti{B}{W\emid S}
 \le\frac{Q_\rho(\mu_0)+Q_\rho(\mu_1)}2,
\end{equation}
then $\muti{B}{S_\rho,W}\le Q_\rho(\mu)$.
\end{lemma}

The coordinate condition need not hold on two inputs: for conjunction,
both mean gaps are $1/2$, while $\nu=3/4$. We therefore start the
induction with two inputs.

\begin{lemma}\label{lem:two-bit}
The bound in \autoref{thm:ck} holds for functions of at most two bits.
\end{lemma}
\begin{proof}
Up to input permutations, input complements, and an output complement,
such a function is a constant, a coordinate, parity, or conjunction.
Constants and coordinates attain $Q_\rho$. For parity, $\mu=1/2$ and
\[
 \muti{B}{X}=\kldivb{(1-\rho^2)/2}{1/2}
 \le\kldivb{p}{1/2}=Q_\rho(1/2).
\]
For conjunction, the four conditional probabilities of $B=1$ given
$X$ are $p^2,p(1-p),p(1-p),(1-p)^2$. Thus the slack
$Q_\rho(1/4)-\muti{B}{X}$ equals
\[
 g(p)=Q_\rho(1/4)-\frac14\left(
 \kldivb{p^2}{1/4}+2\kldivb{p(1-p)}{1/4}
 +\kldivb{(1-p)^2}{1/4}\right).
\]
We have $g(0)=g(1/2)=g'(1/2)=0$. With $w=p(1-p)$,
differentiation gives
\[
 g'''(p)=\frac{(1-2p)(5w^4+2w^3+39w^2+4w+4)}
 {4(\ln2)w^2(2+w)^2(1-w)^2}>0
 \qquad(0<p<1/2).
\]
For $x\in(0,1/2)$, set $c=g(x)/[x(x-1/2)^2]$.
The function $g(t)-ct(t-1/2)^2$ vanishes at $0,x,1/2$
and has derivative zero at $1/2$. Three applications of Rolle's
theorem give $6c=g'''(\xi)>0$ for some $\xi\in(0,1/2)$.
Hence $g(x)>0$, completing the base case.
\end{proof}

\begin{proof}[Proof of \autoref{thm:ck}]
At $\rho=0$, $B$ is independent of $X$ and $Q_0(\mu)=0$.
At $\rho=1$, $Y=X$ and
$\muti{B}{X}=\muti{B}{Y}=Q_1(\mu)$.
Fix $0<\rho<1$. The case $n\le2$ is \autoref{lem:two-bit}.
For $n\ge3$, choose $i$ by \autoref{lem:coordinate-gap} and
write $f_s$ for the restriction of $f$ to $Y_i=s$.
Put $B_s=f_s(Y_{-i})$ and $\mu_s=\E[B_s]$.
The induction hypotheses give
\[
 \muti{B_s}{X_{-i}}\le Q_\rho(\mu_s),\qquad s\in\binary.
\]
Since $Y_i$ is independent of $(X_{-i},Y_{-i})$,
\[
 \muti{B}{X_{-i}\emid Y_i}
 =\frac{\muti{B_0}{X_{-i}}+\muti{B_1}{X_{-i}}}{2}
 \le\frac{Q_\rho(\mu_0)+Q_\rho(\mu_1)}2.
\]
The uniform binary symmetric channel is unchanged when its input
and output are exchanged. Consequently \autoref{lem:closure}
applies with $S=Y_i$, $W=X_{-i}$, and $S_\rho=X_i$.
Our choice of $i$ ensures its mean-gap condition, so it yields
$\muti{B}{X}\le Q_\rho(\mu)$.
\end{proof}

Together with \autoref{lem:q}, this proves \eqref{eq:ck}.
Noise levels $p>1/2$ reduce to $1-p$ by complementing every
coordinate of $Y$.

\subsection{The recombination estimate}
The first ingredient of \autoref{lem:closure} quantifies the dependence
between $S$ and $W$ created by conditioning on $B$.
For $1/2\le\mu<1$, define
\begin{equation}\label{eq:app-allowance}
 a(\mu)=\nu\left(1+\mu-\frac{(1-\mu)^2}{3}\right),
\end{equation}
and set $a(\mu)=a(1-\mu)$ when $0<\mu<1/2$.

\begin{lemma}\label{lem:consistency}
In the setting of \autoref{lem:closure}, assume $0<\mu<1$ and let
$\gamma=(1-\rho^2)d^2/(a(\mu)+d^2)$. Then
\begin{equation}\label{eq:recombination-information}
 \muti{B}{S_\rho,W}
 \le\muti{B}{S_\rho}+(1-\gamma)\muti{B}{W\emid S}.
\end{equation}
The first term depends only on the two means:
\[
 \muti{B}{S_\rho}
 =\frac12\left(\kldivb{\mu-\rho d/2}{\mu}
                  +\kldivb{\mu+\rho d/2}{\mu}\right).
\]
\end{lemma}
\begin{proof}
Complementing $B$ and $S$ as needed, assume $\mu\ge1/2$ and
$\mu_1-\mu_0=d$. Write
\[
 I_0=\muti{B}{W},\qquad T=\muti{S}{W\emid B},\qquad
 J_0=\muti{B}{W\emid S}=I_0+T.
\]
The chain rule and $S\perp W$ give
\begin{equation}\label{eq:app-recombine}
 \muti{B}{S_\rho,W}
 =\muti{B}{S_\rho}+J_0-T+\muti{S_\rho}{W\emid B}.
\end{equation}
Binary divergence contracts under the channel:
\[
 \kldivb{(1+\rho u)/2}{(1+\rho v)/2}
 \le\rho^2\kldivb{(1+u)/2}{(1+v)/2}.
\]
Indeed, the two sides are the Bregman divergences of
$\kappa(\rho t)$ and $\rho^2\kappa(t)$, whose curvatures satisfy
\[
 \frac{\rho^2}{(\ln2)(1-\rho^2t^2)}
 \le\frac{\rho^2}{(\ln2)(1-t^2)}.
\]
Integrating between the means and averaging conditionally on $B$ gives
$\muti{S_\rho}{W\emid B}\le\rho^2T$.
It remains to show $T\ge d^2J_0/(a(\mu)+d^2)$.

The case $d=0$ is immediate. Otherwise put
$c=\Pr(B=1\mid W)$, let $t$ be the difference of the two conditional
probabilities given $(S,W)$, and set
$\sigma^2=\operatorname{Var}(c)$.
Then $\E c=\mu$, $\E t=d$. The posterior sign means of $S$ given
$(B,W)$ are $t/(2c)$ and $-t/[2(1-c)]$. Since
\[
 \kldivb{(1+u)/2}{(1+v)/2}\ge\frac{(u-v)^2}{2\ln2},
\]
averaging and applying weighted Cauchy--Schwarz yields
\[
 T\ge\frac1{8\ln2}\left(\E\frac{t^2}{c(1-c)}-\frac{4d^2}{\nu}\right)
 \ge\frac{2d^2\sigma^2}{(\ln2)\nu(\nu-4\sigma^2)}
 \ge\frac{2d^2\sigma^2}{(\ln2)\nu^2}.
\]
At $c=0,1$, feasibility forces $t=0$; the ratio is defined as zero.
Also $\nu-4\sigma^2=4\E[c(1-c)]>0$, since otherwise $d=0$.

For $\mu\ge1/2$, the complementary estimate is
\begin{equation}\label{eq:binary-kl-upper}
 \kldivb{c}{\mu}
 \le\frac{-\ln\mu}{(\ln2)(1-\mu)^2}(c-\mu)^2.
\end{equation}
To see this, the quotient by $(c-\mu)^2$ is an average of
$1/[(\ln2)x(1-x)]$ along the segment from $\mu$ to $c$, with
weight $1-t$. It is convex in $c$, so its maximum is at an endpoint.
The endpoint $c=1$ is larger: their difference has the sign of
$-\mu^2\ln\mu+(1-\mu)^2\ln(1-\mu)$, a concave function vanishing
at $1/2$ and $1$.
Finally, the logarithm series gives
\[
 \frac{8\mu^2(-\ln\mu)}\nu
 =2\left(1-\sum_{k\ge1}\frac{(1-\mu)^k}{k(k+1)}\right)
 \le1+\mu-\frac{(1-\mu)^2}{3}=\frac{a(\mu)}\nu.
\]
Consequently
\[
 I_0=\E\sparen{\kldivb{c}{\mu}}
 \le\frac{2a(\mu)\sigma^2}{(\ln2)\nu^2}.
\]
The two estimates imply $T\ge d^2I_0/a(\mu)$, hence
$T\ge d^2J_0/(a(\mu)+d^2)$. Substitution into
\eqref{eq:app-recombine} proves the lemma.
\end{proof}
